% !TeX root = ../main.tex
% Fonti lette sulle immagini: Lezione 2 AL.pdf, pp. 2--11;
% Lezione 3.pdf, pp. 2--4. Nessuna pagina di contenuto di questi intervalli omessa.
% La prova del binomio di Newton e le interpretazioni combinatorie sviluppano
% gli enunciati sintetici della fonte. I registri sono in lavoro/SESSIONI.md.
\chapter{Numeri naturali, induzione e combinatoria}
\label{cap:naturali-combinatoria}
\fonti{Lezioni 2 e 3}

\section{I naturali e il principio del minimo intero}
% Fonte: Lezione 2, p.~2

Gli insiemi numerici che incontreremo si dispongono nella catena
\[
  \N\subseteq\Z\subseteq\Q\subseteq\R\subseteq\mathbb C,
\]
dove $\mathbb C$ indica l'insieme dei numeri complessi, richiamato qui
soltanto per collocare i diversi sistemi numerici. Partiamo dai naturali,
che comprendono lo zero:
\[
  \N=\{0,1,2,3,\ldots\}.
\]
L'ordine usuale $\le$ rende $\N$ totalmente ordinato: dati due naturali,
è sempre possibile confrontarli. Per lavorare sui naturali si assume
inoltre la seguente proprietà, più forte del solo ordine totale.

\begin{assioma}[Principio del minimo intero]
Ogni sottoinsieme non vuoto di $\N$ ammette un minimo. In formule,
\[
  S\subseteq\N,\quad S\ne\varnothing
  \quad\Longrightarrow\quad
  \exists m\in S:\ \forall n\in S,\ m\le n.
\]
\end{assioma}

Il minimo deve appartenere a $S$ ed essere minore o uguale a tutti gli
altri suoi elementi. Il principio permette di parlare del \emph{primo}
naturale che soddisfa una certa proprietà; in una dimostrazione per assurdo
permetterà, per esempio, di scegliere il primo eventuale controesempio.

\section{Il principio di induzione}
% Fonte: Lezione 2, pp.~3 e~8

Il principio di induzione formalizza un'idea semplice: se una proprietà
vale all'inizio e passa da ogni naturale al suo successivo, allora vale
per tutti i naturali a partire da quello iniziale. Possiamo esprimerlo
in due forme equivalenti.

\begin{proposizione}[Induzione: forma insiemistica]
Sia $S\subseteq\N$. Se $0\in S$ e, per ogni $n\in S$, anche
$n+1\in S$, allora $S=\N$.
\end{proposizione}

\begin{proposizione}[Induzione: forma con un predicato]
Sia $n_0\in\N$ e sia $P(n)$ un predicato definito per i naturali
$n\ge n_0$. Supponiamo che valgano:
\begin{enumerate}
  \item la \emph{base induttiva}: $P(n_0)$ è vera;
  \item il \emph{passo induttivo}: per ogni $n\ge n_0$,
    $P(n)\Rightarrow P(n+1)$.
\end{enumerate}
Allora $P(n)$ è vera per ogni $n\ge n_0$.
\end{proposizione}

\begin{proof}
Ragioniamo per assurdo.
Supponiamo che esista almeno un naturale $n\ge n_0$ per cui $P(n)$ è falsa.
L'insieme dei controesempi
\[
  F=\{n\in\N:n\ge n_0\text{ e }P(n)\text{ è falsa}\}
\]
è allora non vuoto. Per il principio del minimo intero esiste
$\bar n=\min F$. Non può essere $\bar n=n_0$, perché $P(n_0)$ è vera
per la base induttiva. Dunque $\bar n>n_0$, e il naturale precedente
$\bar n-1$ appartiene ancora all'intervallo di indici considerato.

Poiché $\bar n$ è il primo controesempio, $\bar n-1\notin F$:
$P(\bar n-1)$ è vera. Applicando il passo induttivo con $n=\bar n-1$
otteniamo che anche $P(\bar n)$ è vera, in contraddizione con
$\bar n\in F$. Pertanto $F$ è vuoto e la proprietà vale per tutti
i naturali $n\ge n_0$.
\end{proof}

La forma insiemistica si ricava scegliendo $n_0=0$ e ponendo
$P(n)$ uguale alla proprietà «$n\in S$». Viceversa, l'insieme degli indici
per cui un predicato è vero rende esplicito il collegamento fra le due
formulazioni. È essenziale verificare entrambi i passaggi: il solo passo
induttivo non garantisce che esista un punto da cui iniziare.

\section{La somma dei primi naturali}
% Fonte: Lezione 2, pp.~4--6

Consideriamo la somma
\[
  1+2+\cdots+99+100.
\]
Accoppiando il primo termine con l'ultimo, il secondo con il penultimo
e così via, si ottengono $50$ coppie di somma $101$:
\[
  (1+100)+(2+99)+\cdots+(50+51)
  =50\cdot101=\frac{100\cdot101}{2}=5050.
\]
Questo suggerisce una formula valida per un numero qualsiasi di termini.
La notazione $\sum_{k=1}^{n}k$ abbrevia la somma $1+2+\cdots+n$:
$k$ è l'indice che percorre i termini, mentre $n$ indica dove la somma
si arresta.

\begin{proposizione}
Per ogni $n\in\N^*$ vale
\[
  \sum_{k=1}^{n}k=\frac{n(n+1)}2.
\]
\end{proposizione}
\begin{proof}
Procediamo per induzione a partire da $n_0=1$. La base è vera, perché
$\sum_{k=1}^{1}k=1=1(1+1)/2$.

Per il passo induttivo fissiamo $n\ge1$ e assumiamo l'\emph{ipotesi
induttiva} $\sum_{k=1}^{n}k=n(n+1)/2$. La tesi al passo successivo è
\[
  \sum_{k=1}^{n+1}k=\frac{(n+1)(n+2)}2.
\]
Separiamo l'ultimo addendo e usiamo l'ipotesi soltanto per la somma
dei primi $n$ termini:
\begin{align*}
  \sum_{k=1}^{n+1}k
  &=\left(\sum_{k=1}^{n}k\right)+(n+1)\\
  &=\frac{n(n+1)}2+(n+1)
    =\frac{n^2+3n+2}{2}
    =\frac{(n+1)(n+2)}2.
\end{align*}
Abbiamo dimostrato $P(n)\Rightarrow P(n+1)$; per il principio di
induzione la formula vale per ogni $n\ge1$.
\end{proof}

\section{La disuguaglianza di Bernoulli}
% Fonte: Lezione 2, pp.~6--7

\begin{proposizione}[Disuguaglianza di Bernoulli]
Per ogni $n\in\N$ e ogni reale $x>-1$ si ha
\[
  (1+x)^n\ge1+nx.
\]
\end{proposizione}
\begin{proof}
Fissiamo un qualunque $x>-1$ e procediamo per induzione su $n$.
Per $n=0$, entrambi i membri valgono $1$.

Supponiamo ora che $(1+x)^n\ge1+nx$. Poiché $1+x>0$, possiamo
moltiplicare questa disuguaglianza per $1+x$ senza cambiarne il verso:
\begin{align*}
  (1+x)^{n+1}
  &=(1+x)^n(1+x)\\
  &\ge(1+nx)(1+x)\\
  &=1+(n+1)x+nx^2\\
  &\ge1+(n+1)x.
\end{align*}
L'ultimo passaggio usa $nx^2\ge0$. La disuguaglianza è così dimostrata
per $n+1$, e quindi per ogni $n\in\N$. Poiché $x>-1$ era arbitrario,
l'enunciato vale per tutti i valori di $x$ richiesti.
\end{proof}

La condizione su $x$ entra precisamente nel passaggio in cui si moltiplica
l'ipotesi induttiva per $1+x$: anche quando $x$ è negativo, il fattore
$1+x$ resta positivo.

\section{Permutazioni, disposizioni e combinazioni}
% Fonte: Lezione 2, pp.~9--10

Nel calcolo combinatorio si contano modi diversi di ordinare o scegliere
elementi di un insieme finito. Le formule seguenti riguardano
oggetti \emph{distinti} e scelte \emph{senza ripetizione}: un oggetto
già scelto non viene scelto di nuovo.

\subsection{Permutazioni e fattoriale}

Una \emph{permutazione} di $n$ oggetti è un modo di disporli tutti in fila.
Per il primo posto ci sono $n$ possibilità, per il secondo $n-1$ e così
via fino all'ultimo. Il numero di permutazioni è dunque
\[
  P_n=n!=n(n-1)(n-2)\cdots2\cdot1.
\]
Il simbolo $n!$ si legge \emph{fattoriale di $n$}. Si pone $0!=1$:
anche l'insieme vuoto ha un solo ordinamento, quello senza elementi.
Per esempio, tre oggetti distinti hanno $3!=6$ possibili ordinamenti.

\subsection{Disposizioni: scegliere e ordinare}

Una \emph{disposizione} di $k$ oggetti scelti fra $n$ consiste nel
selezionare $k$ elementi distinti e disporli in fila, tenendo conto
dell'ordine. Per $0\le k\le n$ il numero di disposizioni è
\[
  D_{n,k}=n(n-1)\cdots(n-k+1)=\frac{n!}{(n-k)!}.
\]
Per $k=0$ il prodotto è vuoto e vale $1$; per $k=n$ ritroviamo le
permutazioni. Per esempio, per assegnare un primo e un secondo posto
fra cinque persone ci sono $D_{5,2}=5\cdot4=20$ possibilità.

\subsection{Combinazioni: scegliere senza ordinare}

Una \emph{combinazione} di $k$ oggetti fra $n$ è una scelta nella quale
l'ordine non conta: equivale a un sottoinsieme di $k$ elementi.
Indichiamo il numero di queste scelte con $C_{n,k}$. Ogni scelta di
$k$ elementi produce $k!$ disposizioni, perché i suoi elementi possono
essere ordinati in $k!$ modi. Quindi
\[
  D_{n,k}=C_{n,k}\,k!,\qquad
  C_{n,k}=\frac{n!}{k!(n-k)!}.
\]
Per esempio, scegliendo due persone fra cinque senza assegnare ruoli
si ottiene $C_{5,2}=20/2=10$.

\section{Coefficienti binomiali e loro proprietà}
% Fonte: Lezione 2, pp.~10--11; lezione 3, p.~2

Per $n\in\N$ e $k\in\Z$ definiamo il \emph{coefficiente binomiale}
\[
  \binom nk=
  \begin{cases}
    \dfrac{n!}{k!(n-k)!},&0\le k\le n,\\[.8em]
    0,&k<0\text{ oppure }k>n.
  \end{cases}
\]
Quando $0\le k\le n$, esso coincide con $C_{n,k}$. Il valore zero
fuori da questo intervallo consente di scrivere le identità senza
dover separare continuamente i casi agli estremi.

\begin{proposizione}
Per ogni $n\in\N$ valgono le seguenti proprietà:
\begin{enumerate}
  \item $\displaystyle\binom n0=\binom nn=1$;
  \item $\displaystyle\binom nk=0$ se $k<0$ oppure $k>n$;
  \item $\displaystyle\binom nk=\binom n{n-k}$ per ogni $k\in\Z$;
  \item per ogni $k\in\Z$ vale la \emph{relazione di Pascal}:
    \[
      \binom{n+1}k=\binom nk+\binom n{k-1}.
    \]
\end{enumerate}
\end{proposizione}

\begin{proof}
Le prime due proprietà seguono subito dalla definizione e da $0!=1$.
Per $0\le k\le n$, la simmetria si verifica scambiando i fattoriali
al denominatore:
\[
  \binom n{n-k}
   =\frac{n!}{(n-k)!\,[n-(n-k)]!}
   =\frac{n!}{(n-k)!\,k!}=\binom nk.
\]
Se $k$ è fuori dall'intervallo $[0,n]$, lo è anche $n-k$ e i due
coefficienti sono entrambi zero.

Per la relazione di Pascal consideriamo prima $1\le k\le n$ e
portiamo le due frazioni allo stesso denominatore:
\begin{align*}
  \binom nk+\binom n{k-1}
  &=\frac{n!}{k!(n-k)!}
    +\frac{n!}{(k-1)!(n+1-k)!}\\
  &=\frac{n!\bigl[(n+1-k)+k\bigr]}{k!(n+1-k)!}\\
  &=\frac{(n+1)!}{k!(n+1-k)!}
   =\binom{n+1}k.
\end{align*}
Dividendo l'identità per il fattore positivo $n!/[(k-1)!(n-k)!$, la si riduce a
\[
  \frac{n+1}{k(n+1-k)}
   =\frac1k+\frac1{n+1-k}
   =\frac{n+1-k+k}{k(n+1-k)}.
\]
Per $k=0$ e $k=n+1$ la formula diventa rispettivamente $1=1+0$
e $1=0+1$; per gli altri indici esterni entrambi i membri sono zero.
\end{proof}

La simmetria ha anche un significato combinatorio: scegliere $k$
elementi equivale a scegliere gli $n-k$ da escludere. Analogamente,
per scegliere $k$ elementi fra $n+1$ si può fissare un elemento e
distinguere le scelte che non lo contengono, in numero $\binom nk$,
da quelle che lo contengono, in numero $\binom n{k-1}$.

\section{Il triangolo di Tartaglia}
% Fonte: Lezione 3, p.~3

Scrivendo in ogni riga i coefficienti $\binom n0,\binom n1,\ldots,
\binom nn$ si ottiene il \emph{triangolo di Tartaglia}. Le estremità
valgono $1$ e ciascun elemento interno è la somma dei due che lo
precedono nella riga superiore: è proprio la relazione di Pascal.

\begin{schema}
\begin{tikzpicture}[x=1.05cm,y=.85cm]
  \foreach \riga/\valori in {0/{1},1/{1,1},2/{1,2,1},3/{1,3,3,1},4/{1,4,6,4,1}} {
    \foreach \valore [count=\colonna from 0] in \valori {
      \node (t-\riga-\colonna) at ({\colonna-\riga/2},-\riga) {$\valore$};
    }
    \node[anchor=east,black!65] at (-3,-\riga) {$n=\riga$};
  }
  \draw[coloreschema,-{Stealth[length=1.8mm]},shorten <=3pt,shorten >=3pt]
    (t-3-1) -- (t-4-2);
  \draw[coloreschema,-{Stealth[length=1.8mm]},shorten <=3pt,shorten >=3pt]
    (t-3-2) -- (t-4-2);
  \node[coloreschema,anchor=west] at (2.9,-3.7) {$3+3=6$};
\end{tikzpicture}
\par\smallskip
Ogni riga contiene i coefficienti binomiali di un dato ordine $n$.
Le frecce mostrano come due coefficienti della riga $n=3$ generano
quello centrale della riga $n=4$.
\end{schema}

\section{La formula del binomio di Newton}
% Fonte: Lezione 3, p.~4

\begin{proposizione}[Binomio di Newton]
Per $a,b\in\R$ e $n\in\N$ vale
\[
  (a+b)^n=\sum_{k=0}^{n}\binom nk a^{n-k}b^k.
\]
I fattori con esponente zero si interpretano come $1$, secondo
la convenzione usata nelle identità polinomiali.
\end{proposizione}

Il termine di indice $k$ contiene $k$ fattori $b$ e $n-k$ fattori $a$.
Per ottenerlo dal prodotto di $n$ parentesi $(a+b)$ si scelgono le
$k$ parentesi dalle quali prendere $b$: questo spiega la presenza
del coefficiente $\binom nk$.

\begin{proof}
Procediamo per induzione. Per
$n=0$ la formula vale, perché entrambi i membri sono $1$ con la
convenzione appena precisata. Supponiamo ora che valga per $n$.
Moltiplicando per $a+b$ e distribuendo i due fattori otteniamo
\begin{align*}
  (a+b)^{n+1}
  &=\sum_{k=0}^{n}\binom nk a^{n+1-k}b^k
    +\sum_{k=0}^{n}\binom nk a^{n-k}b^{k+1}\\
  &=\sum_{k=0}^{n}\binom nk a^{n+1-k}b^k
    +\sum_{k=1}^{n+1}\binom n{k-1}a^{n+1-k}b^k.
\end{align*}
Nella seconda somma abbiamo cambiato indice, in modo che le potenze
di $a$ e $b$ abbiano la stessa forma. Per gli indici $1\le k\le n$
il coefficiente complessivo è allora
$\binom nk+\binom n{k-1}=\binom{n+1}k$.
Ai due estremi restano $a^{n+1}$ e $b^{n+1}$, entrambi con coefficiente
$1$. Si ricompone quindi esattamente
$\sum_{k=0}^{n+1}\binom{n+1}k a^{n+1-k}b^k$, che è la formula al passo
successivo. Il principio di induzione conclude la dimostrazione.
\end{proof}

Per esempio, dalla riga $n=4$ del triangolo si legge subito
\[
  (a+b)^4=a^4+4a^3b+6a^2b^2+4ab^3+b^4.
\]
